\documentclass[11pt]{article}
\usepackage[margin=0.95in]{geometry}
\usepackage{amsmath,amssymb}
\usepackage{graphicx}
\usepackage{booktabs,longtable,array}
\usepackage{xcolor}
\usepackage[hidelinks]{hyperref}
\usepackage{enumitem}
\usepackage{caption}
\setlist{itemsep=2pt,topsep=4pt}
\captionsetup{font=small,labelfont=bf}
\setlength{\parskip}{5pt}
\setlength{\parindent}{0pt}
\renewcommand{\arraystretch}{1.15}
\title{Phase 3 Rationale: Removing the Confounds}
\author{}
\date{trm-delay-robust-scaling $\cdot$ phase rationale $\cdot$ 2026-09-14 to 09-23 $\cdot$ 231 jobs\\[2pt]Why the phase existed, what inspired its design, and what it contributed. Results are in \texttt{.\allowbreak{}.\allowbreak{}/\allowbreak{}phase\_\allowbreak{}summaries/\allowbreak{}phase3\_\allowbreak{}summary.\allowbreak{}md}; labels are in \texttt{.\allowbreak{}.\allowbreak{}/\allowbreak{}glossary.\allowbreak{}md}.\\[2pt]Research question (the whole project).** Does weight-tied recursion (TRM), used as the actor-critic trunk of a model-free RL agent, benefit continuous control under action delay, and if not, why not?}
\begin{document}
\maketitle
\section*{In one paragraph}



After Phase 2 the honest position was not "recursion does not work" but "the experiment that would tell us has not been run". The Pendulum deficit could be capacity, the latent reset, or a task with nothing to infer. Phase 3 was designed to remove all three, and to give recursion its best fair chance: a hypothesis (H3) under which TRM \emph{should} win, with a prediction no capacity story could fake. It also added the measurements a control engineer would ask for. The result was a confound-controlled negative: the deficit is architectural, memory and depth recover nothing, and TRM's one advantage, smoother torque, comes with worse robustness and 65$\times$ the latency.



\section*{1. Where the project stood going in}



\begin{itemize}
\item A consistent Pendulum deficit (0/6) with a value-fit deficit (7/7 cells, p = 0.008).
\item Three confounds: C-1 capacity, C-2 latent reset, C-3 nothing to infer.
\item Preconditions P1--P3 for TRM's in-domain success, which the port lacks.
\item 102 banked checkpoints from Phases 1--2 never evaluated as controllers.
\item Phases 1--2 had committed 102 jobs before discovering that most could not answer the question.
\end{itemize}



\section*{2. Why this phase was made}



\begin{enumerate}
\item \textbf{To attribute the deficit.} A deficit that might be parameter count says nothing about recursion. Matching capacity in both directions settles it.
\item \textbf{To give recursion its strongest case.} If TRM is a learned state estimator, it should help where (i) the input leaves something to infer and (ii) the trunk carries information across steps. Phases 1--2 satisfied neither. H3 states that, and its failure would leave only H\_null.
\item \textbf{To judge TRM as a controller, not just a return.} A controls reader needs smoothness, robustness margins and real-time feasibility. None had been measured.
\item \textbf{To spend compute only when it could answer something.} Gated stages, each cheap enough to stop the phase.
\end{enumerate}



\section*{3. Design decisions and what inspired them}



\begin{center}\small\begin{tabular}{>{\raggedright\arraybackslash}p{0.172\linewidth}>{\raggedright\arraybackslash}p{0.331\linewidth}>{\raggedright\arraybackslash}p{0.403\linewidth}}
\toprule
Decision & Inspiration & Why it mattered \\
\midrule
H3 with an \emph{interaction} as its load-bearing prediction (H3.2) & Standard logic of confound control: a capacity story cannot flip the sign of a contrast between identical networks across conditions & Carry-Z $-$ MLP $\leq$ 0 with nothing to infer, > 0 with something to infer. A main effect could be capacity; a sign flip could not. \\
Competing accounts H\_cap and H\_null stated up front & Pre-registration practice & Every outcome had a named reading before the data. \\
Six gated stages with rules written before compute & Phases 1--2's sunk 102 jobs & Each stage could stop the phase; Stages 3--4 were in fact gated off when G2 failed. \\
Headroom first (S0.1, G0) & "Before asking whether recursion exploits uncertainty, check that uncertainty costs anything" & Hiding the queue cost the MLP 70.8 (6/6), so memory had something to recover. \\
The 2$\times$2 capacity design plus H = 1 and H = 6 (Stage 1) & Growing the MLP alone only moves the comparison; shrinking TRM as well allows attribution & G1: TRM loses at 136k and 788k; MLP capacity inert. Depth inert too, since H = 1 removes the outer recursion with nothing else changed. \\
Carry-Z with stored-state replay & HRM-Agent (Dang \& Rawlinson, 2025) for carrying the latent; R2D2 (Kapturowski et al., 2019) for replaying it & Carrying the latent at acting time but resetting it at update time would silently confound the test. Replay was verified to reproduce the behaviour policy exactly. \\
A carried LSTM arm & The standard recurrent POMDP baseline; Ni et al. (2022) & Separates "memory helps" from "TRM's memory helps". \\
Previous-action input for all trunks (A-6) & Ni et al. (2022) & Trunks then differ only in architecture and memory. \\
Amendment A-1: hide the queue (C4) & A theory check before submission: with a delay redrawn every step, the queue already contains every action that can land & An oracle arm could only be null; hiding the queue creates something genuine to infer. \\
Latent per-episode delay (C6, G2s) & An independent plan review (2026-09-17), read after execution began & A delay drawn once per episode is something memory can identify. Run after G2 failed, with its rule fixed before its data, labelled supplementary. \\
Control-quality battery and latency benchmark (S0.2) & Control-engineering practice: a controller is judged on its envelope, not its return & Measured torque reversals, robustness to noise, mass and dropout, and per-decision latency against control periods. \\
IQM and stratified bootstrap alongside sign tests & Agarwal et al. (2021) & Reporting standards for few-seed RL. \\
Frozen code snapshots, ledgers, verification before data & Lessons from Phase 2's filename collision & Every result traceable to exact code. Carried into Phases 4 and 5. \\
\bottomrule
\end{tabular}\end{center}



\section*{4. How it contributed to the research question}



\begin{center}\small\begin{tabular}{>{\raggedright\arraybackslash}p{0.164\linewidth}>{\raggedright\arraybackslash}p{0.493\linewidth}>{\raggedright\arraybackslash}p{0.249\linewidth}}
\toprule
Hypothesis & Outcome & What it means for the question \\
\midrule
H\_cap, the deficit is capacity (G1) & \textbf{Rejected.} TRM 136k $-$ MLP 135k $-$50.2 (0/12); TRM 788k $-$ MLP 788k $-$27.1 (2/12); MLP 788k $-$ 135k $-$3.8 (5/12) & The deficit is architectural. \\
H3.2, uncertainty $\times$ memory (G2) & \textbf{Rejected.} Wrong sign at C4 ($-$40.0, 1/6) & Recursion does not exploit residual uncertainty here. \\
Carrying the latent (G3, G2s) & \textbf{Fail.} Carry-Z never beats Reset-Z in $\geq$ 5/6 & Confound C-2 removed; nothing recovered. \\
Depth & Inert: H = 1 and H = 6 vs H = 3 give $-$3.3 (5/12) and $-$10.9 (4/12) & Extra recursion buys nothing. \\
Value fit & TRM's critic fits worse at matched capacity (1.45$\times$, 1.89$\times$, 12/12 each) & The mechanism of the deficit. \\
H3.4 smoothness & \textbf{Supported}: 0.4 vs 9 torque reversals/s (11/12, p = 0.006) & TRM's one advantage. \\
H3.5 graceful degradation & \textbf{Contradicted}: worse under sensor noise, a heavier plant, dropout &  \\
H3.6 real-time feasibility & 4.13 vs 0.064 ms p99 per decision (65$\times$); 100 Hz feasible, 500 Hz not & A deployment cost. \\
\bottomrule
\end{tabular}\end{center}



Both halves of the question received their main answer here. \textbf{Does it benefit?} No, with capacity matched, the latent carried, and conditions built to leave something to infer. \textbf{Why not?} Not capacity, not memory, not depth; the deficit sits in value fitting and mostly in training (the final greedy policies differ by only $-$3.6, 2/12).



\section*{5. How it shaped the experiments that followed}



Phase 3's gates kept its confirmatory claims clean, which left five questions open. Phase 4 took all five under its own pre-registration:



\begin{center}\small\begin{tabular}{>{\raggedright\arraybackslash}p{0.517\linewidth}>{\raggedright\arraybackslash}p{0.417\linewidth}}
\toprule
Left open by Phase 3 & Phase 4 stage \\
\midrule
TRM's training discipline (full gradient, per-cycle supervision) & Stage A (GA) \\
Was TRM ever tuned? & Stage C (GC) \\
Is the attribution DDPG-specific? & Stage T (GT), then Phase 5 \\
An unstable plant, and complexity & Stage B0 on InvertedDoublePendulum, Hopper fallback \\
Why is TRM smoother yet less robust? & Stage E (Lipschitz analysis) \\
\bottomrule
\end{tabular}\end{center}



\section*{6. In hindsight}



\begin{itemize}
\item \textbf{The plan review should have been read before execution.} Execution reached its central point independently (A-1) but not its recommended fix, which then had to run as a supplementary test.
\item \textbf{The interaction design worked as intended:} it gave a result no capacity argument can explain away.
\item \textbf{The battery changed the story's texture.} "TRM loses on return" became "TRM is a smoother, slower, more fragile controller", which is what an engineer needs to know.
\end{itemize}



\section*{Sources}



\begin{itemize}
\item Dang \& Rawlinson (2025), HRM-Agent; Kapturowski et al. (2019), R2D2; Ni et al. (2022), recurrent model-free RL; Agarwal et al. (2021), RL evaluation statistics; Derman et al. (2021), non-stationary delays (planned C5, not run).
\item Project files:
\begin{itemize}
\item \texttt{phase3/\allowbreak{}plan/\allowbreak{}plan.\allowbreak{}md} (\S{}\S{}2--4, \S{}6 exploration inventory, \S{}10 amendments)
\item \texttt{phase3/\allowbreak{}report/\allowbreak{}}
\item \texttt{.\allowbreak{}.\allowbreak{}/\allowbreak{}full\_\allowbreak{}report.\allowbreak{}md} \S{}\S{}8--9
\end{itemize}
\end{itemize}
\end{document}
